%----------------------------------------------------------------------------------------
%	DOCUMENT VARIABLES
%	Fill in the lines below to enter your information into the thesis template
%	Each of the commands can be cited anywhere in the thesis
%----------------------------------------------------------------------------------------

\newcommand{\myTitle}{Smart gateway for Low-Power Lossy Networks\xspace}
\newcommand{\mySubtitle}{A PhD thesis\xspace}
\DeclareRobustCommand{\myName}{Remy Leone\xspace}

\newcommand{\Paolo}{Dr. Paolo MEDAGLIANI\xspace}
\newcommand{\Jeremie}{Dr. Jeremie LEGUAY\xspace}
\newcommand{\Claude}{Dr. Claude CHAUDET\xspace}
\newcommand{\Vania}{Dr. Vania CONAN\xspace}
\newcommand{\JeanLouis}{Dr. Jean Louis ROUGIER\xspace}

\newcommand{\myFaculty}{INFRES\xspace}
\newcommand{\myDepartment}{Network and Computer Science Department\xspace}
\newcommand{\myUni}{Telecom Paris Tech\xspace}
\newcommand{\myLocation}{Paris, France\xspace}
\newcommand{\myTime}{\today\xspace}

%----------------------------------------------------------------------------------------
%	USEFUL COMMANDS
%----------------------------------------------------------------------------------------

\newcommand{\ie}{i.\,e.}
\newcommand{\Ie}{I.\,e.}
\newcommand{\eg}{e.\,g.}
\newcommand{\Eg}{E.\,g.} 

\newcounter{dummy} % Necessary for correct hyperlinks (to index, bib, etc.)
% \providecommand{\mLyX}{L\kern-.1667em\lower.25em\hbox{Y}\kern-.125emX\@}

%----------------------------------------------------------------------------------------
%	PACKAGES
%----------------------------------------------------------------------------------------

%\usepackage[francais]{babel}
\usepackage{lipsum} % Used for inserting dummy 'Lorem ipsum' text into the template 
\usepackage[utf8]{inputenc} 
\usepackage[fleqn]{amsmath} % Math environments and more by the AMS 
\usepackage{xspace} % To get the spacing after macros right
\usepackage{mparhack} % To get marginpar right
\usepackage{fixltx2e} % Fixes some LaTeX stuff 

\PassOptionsToPackage{smaller}{acronym} % Include printonlyused in the first bracket to only show acronyms used in the text
\usepackage{acronym} % nice macros for handling all acronyms in the thesis

%------------------------------------------------

%\renewcommand*{\acsfont}[1]{\textssc{#1}} % For MinionPro
\renewcommand{\bflabel}[1]{{#1}\hfill} % Fix the list of acronyms


% \usepackage[backend=biber,style=alphabetic,sorting=ynt]{biblatex}
% \addbibresource{main.bib}
\usepackage[pdftex]{graphicx} 
\usepackage{epstopdf}

\usepackage[l2tabu, orthodox]{nag} % Static Analysis of LaTeX source code

%------------------------------------------------
% Telecom Paris Tech
%------------------------------------------------

\usepackage{geometry}
\usepackage{eso-pic}
\usepackage{ifpdf}

\newcommand\BackgroundPicLastPage{
\ifpdf
	\includegraphics[height=\paperheight,width=\paperwidth]{img/logos/cover_4_bg.pdf}
\else
	\includegraphics[height=\paperheight,width=\paperwidth]{img/logos/cover_4_bg.pdf}
\fi
}

\newcommand\BackgroundPicCover{
\ifpdf
	\includegraphics[height=\paperheight,width=\paperwidth]{img/logos/cover_bg.pdf}
\else
	\includegraphics[height=\paperheight,width=\paperwidth]{img/logos/cover_bg.pdf}
\fi}

%----------------------------------------------------------------------------------------
%	FLOATS: TABLES, FIGURES AND CAPTIONS SETUP
%----------------------------------------------------------------------------------------

\usepackage{tabularx} % Better tables
\setlength{\extrarowheight}{3pt} % Increase table row height
\newcommand{\tableheadline}[1]{\multicolumn{1}{c}{\spacedlowsmallcaps{#1}}}
\newcommand{\myfloatalign}{\centering} % To be used with each float for alignment
\usepackage{caption}
\captionsetup{format=hang,font=small}
\usepackage{subfig}  

%----------------------------------------------------------------------------------------
%	CODE LISTINGS SETUP
%----------------------------------------------------------------------------------------

\usepackage{listings} 
%\lstset{emph={trueIndex,root},emphstyle=\color{BlueViolet}}%\underbar} % for special keywords
\lstset{language=[LaTeX]Tex, % Specify the language for listings here
% keywordstyle=\color{Blue}, % Add \bfseries for bold
basicstyle=\small\ttfamily, % Makes listings a smaller font size and a different font
%identifierstyle=\color{NavyBlue}, % Color of text inside brackets
commentstyle=\color{Green}\ttfamily, % Color of comments
stringstyle=\rmfamily, % Font type to use for strings
numbers=left, % Change left to none to remove line numbers
numberstyle=\scriptsize, % Font size of the line numbers
stepnumber=5, % Increment of line numbers
numbersep=8pt, % Distance of line numbers from code listing
showstringspaces=false, % Sets whether spaces in strings should appear underlined
breaklines=true, % Force the code to stay in the confines of the listing box
%frameround=ftff, % Uncomment for rounded frame
frame=single, % Frame border - none/leftline/topline/bottomline/lines/single/shadowbox/L
belowcaptionskip=.75\baselineskip % Space after the "Listing #: Desciption" text and the listing box
}

%----------------------------------------------------------------------------------------
%	HYPERREFERENCES
%----------------------------------------------------------------------------------------

\PassOptionsToPackage{pdftex,hyperfootnotes=false,pdfpagelabels}{hyperref}
\usepackage{hyperref}  % backref linktocpage pagebackref
\pdfcompresslevel=9
\pdfadjustspacing=1

\hypersetup{
% Uncomment the line below to remove all links (to references, figures, tables, etc)
%draft, 
colorlinks=true, linktocpage=true, pdfstartpage=3, pdfstartview=FitV,
% Uncomment the line below if you want to have black links (e.g. for printing black and white)
%colorlinks=false, linktocpage=false, pdfborder={0 0 0}, pdfstartpage=3, pdfstartview=FitV, 
breaklinks=true, pdfpagemode=UseNone, pageanchor=true, pdfpagemode=UseOutlines,
plainpages=false, bookmarksnumbered, bookmarksopen=true, bookmarksopenlevel=1,
hypertexnames=true, pdfhighlight=/O, urlcolor=webbrown, %linkcolor=RoyalBlue, citecolor=webgreen,
%------------------------------------------------
% PDF file meta-information
pdftitle={\myTitle},
pdfauthor={\textcopyright\ \myName, \myUni, \myFaculty},
pdfsubject={},
pdfkeywords={},
pdfcreator={pdfLaTeX},
pdfproducer={LaTeX with hyperref}
%------------------------------------------------
}   

%----------------------------------------------------------------------------------------
%	BACKREFERENCES
%----------------------------------------------------------------------------------------

% \usepackage{ifthen} % Allows the user of the \ifthenelse command
% \newboolean{enable-backrefs} % Variable to enable backrefs in the bibliography
% \setboolean{enable-backrefs}{true} % Variable value: true or false

% \newcommand{\backrefnotcitedstring}{\relax} % (Not cited.)
% \newcommand{\backrefcitedsinglestring}[1]{(Cited on page~#1.)}
% \newcommand{\backrefcitedmultistring}[1]{(Cited on pages~#1.)}
% \ifthenelse{\boolean{enable-backrefs}} % If backrefs were enabled
% {
% \PassOptionsToPackage{hyperpageref}{backref}
% \usepackage{backref} % to be loaded after hyperref package 
% \renewcommand{\backreftwosep}{ and~} % separate 2 pages
% \renewcommand{\backreflastsep}{, and~} % separate last of longer list
% \renewcommand*{\backref}[1]{}  % disable standard
% \renewcommand*{\backrefalt}[4]{% detailed backref
% \ifcase #1 
% \backrefnotcitedstring
% \or
% \backrefcitedsinglestring{#2}
% \else
% \backrefcitedmultistring{#2}
% \fi}
% }{\relax} 

%----------------------------------------------------------------------------------------
%	CHANGING TEXT AREA 
%----------------------------------------------------------------------------------------

\usepackage[T1]{fontenc}
\usepackage[bitstream-charter]{mathdesign} % Font
\usepackage[Lenny]{fncychap} % Nice chapter headings Other choices : Sonny, Lenny, Glenn, Conny, Rejne, Bjarne

\usepackage{minitoc} % Mini table of contents
\usepackage{epigraph} % This is to insert the nice quote at the beginning of the chapter 

%------------------------------------------------
% Page headers fancy style  
%------------------------------------------------

\usepackage{fancyhdr}